\snote{identif}

\subsection*{Question}
The method selects a subset $S$ that is \emph{sufficient}: re-estimating $S$ together with the calibration parameters $\Theta_C$, with the other parameters at $\theta_0$, reproduces the outputs of the full model. Sufficiency does not imply that the data determine the values of the selected parameters. A selected parameter could, for example, compensate for a scaling factor of the measurement, so that only their product is determined. This note examines whether the selected parameters are also practically identifiable in the reduced problem, that is, with $S\cup\Theta_C$ free and the remaining parameters at their published values.

\subsection*{Design}
\begin{enumerate}\itemsep0pt
\item \emph{Local check.} At $\theta_0$, the columns of the sensitivity matrix for $S\cup\Theta_C$, normalised to unit length and not centred (the Fisher information of the reduced problem up to scaling), were examined: for each selected parameter, the variance inflation factor (VIF) within $S\cup\Theta_C$ and the column of $\Theta_C$ or $S$ most similar to it.
\item \emph{Profile likelihood of the reduced problem.} With all measurements of the problem and the noise levels of the benchmark, the reduced problem was fitted from $\theta_0$ (minimum $\hat\theta$). For each selected parameter $\theta_j$, $\log|\theta_j|$ was fixed at displacements of 0.05 to 4.6 from $\log|\hat\theta_j|$ and the other parameters of $S\cup\Theta_C$ were re-estimated, with the same least-squares method and bounds $[|\theta_0|/10,10|\theta_0|]$ as in the rest of the study. If a profile point improved the minimum, the reduced problem was re-fitted from that point and the profile repeated (at most three rounds). Only points within the bounds of $\theta_j$ were used. As in Raue et al.\ (2009), $\theta_j$ is \emph{identifiable} if $\Delta\chi^2$ exceeds 3.84 (95\%, one degree of freedom) on both sides within the bounds, \emph{practically non-identifiable} if on one side only, and \emph{flat} otherwise. Because several problems have recorded noise levels that do not match the scatter of the data, a robust variant scales the threshold by $\max(1,\chi^2_{\min}/\mathrm{d.f.})$, which is equivalent to rescaling the noise so that the fit has $\chi^2/\mathrm{d.f.}=1$.
\item \emph{Condition of the formal theorem.} For each selected parameter we computed the quantities that appear in the theorem described below: the profiled information $I_j=1/[(A^{\top}A)^{-1}]_{jj}$, where $A$ is the sensitivity matrix of the residuals with respect to $\log|\theta|$ for $S\cup\Theta_C$ at $\hat\theta$; the residual gradient $\gamma=\|P r_0\|$, the norm of the projection of the residual at $\hat\theta$ onto the range of $A$; and the nonlinear remainder $\|r(\hat\theta+\delta)-r_0-A\delta\|$ at the computed profile points.
\end{enumerate}
Six problems with expensive simulations (Bachmann, Chen, Isensee, Lang, Raimundez, Smith) were computed on a 256-core server and the others on a workstation.

\subsection*{Formal result}
The Lean module \texttt{ProfileIdentifiability} (Supplementary Table~\ref{tab:statements}) proves, for any residual function $r$, linear map $A$ and constants, the following statement. Let $\Delta\chi^2(\delta)=\|r(\hat\theta+\delta)\|^2-\|r_0\|^2$. If $|\langle r_0,A\delta\rangle|\le\gamma\|A\delta\|$ for all $\delta$ (approximate stationarity of $\hat\theta$), $I_j\,\delta_j^2\le\|A\delta\|^2$ for all $\delta$ (which holds with $I_j=1/[(A^{\top}A)^{-1}]_{jj}$), and the nonlinear remainder is at most $\varepsilon$ on the admissible set, then every admissible $\delta$ with
\[ |\delta_j|\ \ge\ t^{\ast}=\frac{\varepsilon+\gamma+\sqrt{\gamma^2+c+2\|r_0\|\varepsilon+2\gamma\varepsilon}}{\sqrt{I_j}} \]
has $\Delta\chi^2(\delta)\ge c$. Hence, if $t^{\ast}$ is smaller than the distance from $\log|\hat\theta_j|$ to the bounds on both sides, $\theta_j$ is practically identifiable within the bounds (\texttt{practically\_identifiable\_grad}). For a linear model ($\varepsilon=0$) and an exact minimum ($\gamma=0$), $t^{\ast}=\sqrt{c\,[(A^{\top}A)^{-1}]_{jj}}$ is the half-width of the Wald interval (\texttt{tstar\_linear}). The module also proves local structural identifiability: if $\|A\delta\|\ge\sigma\|\delta\|$ and the remainder is at most $M\|\delta\|^2$, distinct parameters with $\|\delta\|<\sigma/M$ give distinct outputs (\texttt{local\_injective}).

The theorem is general; its hypotheses are numerical properties of each problem. The stationarity constant $\gamma$ and the information $I_j$ are computed directly. The remainder $\varepsilon$, however, can only be evaluated at points, not bounded over a region, for the same reason that the finite-perturbation certificates of steps~4 and~20 could not be applied (\sref{seq}): far from $\hat\theta$ the response of an ODE model saturates and the linearisation overestimates it. We therefore report two numerical checks of the hypotheses: the linear condition (Wald interval within the bounds, $\varepsilon=0$) and the condition with $\varepsilon$ evaluated at the computed profile points up to the first point at which $|\delta_j|\ge t^{\ast}$.

\subsection*{Results}
Supplementary Table~\ref{tab:identif} gives the results for the 32 problems with an admissible subset.

\emph{Profile likelihood.} In 26 of the 32 problems all selected parameters were practically identifiable, and 67 of the 74 selected parameters were identifiable. The exceptions are of three kinds. In Fujita a selected parameter is confounded with a measurement scaling factor: the columns of \texttt{reaction\_6\_k1} and \texttt{scaling\_pS6\_tot} have cosine 1.00, and neither selected parameter is identifiable. In Giordano and Perelson the recorded noise levels are far larger than the residuals ($\chi^2/n<10^{-2}$), so no parameter change raises $\chi^2$ by 3.84; in Okuonghae and Oliveira one parameter is identifiable and the profile of the other could not be simulated. In Smith, the optimum of \texttt{k8} lies at the lower bound of the study. With the robust threshold, 24 of the 32 problems keep all their parameters identifiable: Raia ($\chi^2/\mathrm{d.f.}=1{,}791$), Chen (168) and Smith (1{,}894) lose one or both, because their recorded noise levels are too small for the scatter of the data.

\emph{Local check.} A VIF above 10 within $S\cup\Theta_C$ flags a correlation between a selected parameter and a calibration parameter, but it does not by itself make the parameter non-identifiable: in Brannmark, Fiedler, Isensee, Lang, Rahman, Raia, Raimundez, Schwen and Weber the VIF exceeds 10, yet all their selected parameters are identifiable in the profiles, because the data are informative enough to separate them. Only in Fujita does the confounding reach the profile.

\emph{Condition of the theorem.} The linear condition agrees with the profile classification for every parameter whose profile could be evaluated (Supplementary Table~\ref{tab:identif}, column Wald), except that a singular $A^{\top}A$ (Schwen, a calibration parameter without effect on the outputs) prevents its evaluation. With the nonlinear remainder evaluated at the profile points, the hypotheses of the theorem hold for a smaller number of parameters (column Lean), and never for a parameter whose profile is not identifiable. The quantitative certificate is therefore conservative: the profiled information $I_j$ explains which parameters are identifiable, and the nonlinearity of the models is what prevents a certificate over the whole range of the study.

\subsection*{Scope}
The result is practical identifiability of the reduced problem: local, within the bounds of the study, with the non-selected parameters at $\theta_0$ and with the noise model of the benchmark. It is not identifiability of the full model, in which all parameters are free: in the profile likelihood of the full model only 20 of the 52 selected parameters evaluated were identifiable (\sref{profile}). That contrast is consistent with the purpose of the method. The data do not determine all the parameters of these models, but, with the other parameters at their published values, a few selected parameters are both sufficient to reproduce the outputs and, in most problems, determined by the data. Identifiability does not imply that the estimated values are the true biological values: they depend on the values at which the other parameters are fixed.

{\small
\begin{longtable}{@{}lrrrrllrr@{}}
\caption{Identifiability of the selected parameters in the reduced problem ($S\cup\Theta_C$ free, other parameters at $\theta_0$, all measurements). $|S|$, $|\Theta_C|$: selected biological and calibration parameters; $\chi^2/n$: at the optimum of the reduced problem; VIF: largest variance inflation factor of a selected parameter among the columns of $J_{S\cup\Theta_C}$ (unit-norm, uncentred) at $\theta_0$. Profile: class of each selected parameter by profile likelihood with $\Delta\chi^2=3.84$ within $[|\theta_0|/10,10|\theta_0|]$ (I, identifiable, the profile exceeds the threshold on both sides; N, practically non-identifiable, on one side; F, flat; --, the profile could not be simulated); robust: the same with the threshold scaled by $\max(1,\chi^2_{\min}/\mathrm{d.f.})$. Wald: parameters whose linear (Wald) interval $\sqrt{3.84\,[(J^{\top}J)^{-1}]_{jj}}$ fits within the bounds on both sides; Lean: parameters for which the hypotheses of \texttt{practically\_identifiable\_grad} hold at the computed profile points up to the first crossing.}\label{tab:identif}\\
\toprule
Problem & $|S|$ & $|\Theta_C|$ & $\chi^2/n$ & VIF & Profile & Robust & Wald & Lean \\
\midrule
\endfirsthead
\toprule
Problem & $|S|$ & $|\Theta_C|$ & $\chi^2/n$ & VIF & Profile & Robust & Wald & Lean \\
\midrule
\endhead
Alkan & 2 & 10 & 1.96 & 1.1 & I I & I I & 2/2 & 1/2 \\
Armistead & 2 & 0 & 1.03 & 1.0 & I I & I I & 2/2 & 1/2 \\
Bachmann & 2 & 74 & 1.16 & 6.1 & F F & F F & -- & -- \\
Bertozzi & 2 & 0 & 1.10 & 5.9 & I I & I I & 2/2 & 0/2 \\
Blasi & 2 & 0 & 6.46 & 2.3 & I I & I I & 2/2 & 2/2 \\
Boehm & 2 & 0 & 1.04 & 1.4 & I I & I I & 2/2 & 2/2 \\
Borghans & 2 & 2 & 0.97 & 2.6 & I I & I I & 2/2 & 0/2 \\
Brannmark & 2 & 3 & 1.13 & 127.2 & I I & I I & 2/2 & 0/2 \\
Bruno & 5 & 0 & 0.90 & 1.3 & I I I I I & I I I I I & 5/5 & 5/5 \\
Elowitz & 2 & 2 & 1.07 & 6.0 & I I & I I & 2/2 & 0/2 \\
Fiedler & 2 & 8 & 1.04 & 47.2 & I I & I I & 2/2 & 0/2 \\
Fujita & 2 & 3 & 5.48 & 2.5e+08 & F F & F F & 0/2 & 0/2 \\
Giordano & 2 & 0 & 0.00 & 4.1 & N F & N F & 0/2 & 0/2 \\
Isensee & 2 & 11 & 1.02 & 41.6 & N F & N F & -- & -- \\
Lang & 2 & 16 & 6.19 & 21.0 & F F & F F & -- & -- \\
Laske & 2 & 2 & 0.18 & 2.7 & I I & I I & 2/2 & 1/2 \\
Liu & 2 & 0 & 1.02 & 1.3 & I I & I I & 2/2 & 0/2 \\
Lucarelli & 10 & 0 & 1.22 & 9.9 & I I I I I I I I I I & I I I I I I I I I I & 10/10 & 4/10 \\
Okuonghae & 2 & 0 & 1,413 & 1.3 & -- I & -- N & 2/2 & 0/2 \\
Oliveira & 2 & 0 & 182 & 1.0 & -- I & -- I & 2/2 & 0/2 \\
Perelson & 2 & 0 & 0.00 & 1.2 & F F & F F & 0/2 & 0/2 \\
Rahman & 2 & 0 & 0.00 & 39.0 & I I & I I & 2/2 & 2/2 \\
Raia & 2 & 5 & 1,791 & 2.3e+04 & I I & F F & 2/2 & 0/2 \\
Raimundez & 2 & 79 & 4.85 & 66.2 & F F & F F & -- & -- \\
Salazar-Cavazos & 2 & 0 & 44.96 & 5.7 & I I & I I & 2/2 & 0/2 \\
Schwen & 2 & 15 & 0.96 & 39.1 & I I & I I & 0/2 & 0/2 \\
Smith & 2 & 9 & 1,894 & 29.3 & F F & F F & -- & -- \\
Sneyd & 2 & 0 & 1.01 & 1.0 & I I & I I & 2/2 & 2/2 \\
Weber & 2 & 5 & 0.94 & 631.0 & I I & I I & 2/2 & 0/2 \\
Zhao & 3 & 0 & 0.27 & 1.0 & I I I & I I I & 3/3 & 1/3 \\
Zheng & 2 & 0 & 1.03 & 1.1 & I I & I I & 2/2 & 2/2 \\
\bottomrule
\end{longtable}

}
